\documentclass[11pt,a4paper]{article}

\usepackage{amsmath,amssymb}
\usepackage{graphicx}
\usepackage{hyperref}
\usepackage{booktabs}
\usepackage[margin=1in]{geometry}
\usepackage{cite}

\title{\textbf{Vehicle License Plate Recognition and Hotlist Alert System for Urban Surveillance}}
\author{Basit Ali \\
Abdul Wali Khan University Mardan, Pakistan \\
\texttt{whoisbasit@gmail.com}}
\date{2025}

\begin{document}
\maketitle

\begin{abstract}
We present a license plate recognition system that detects and reads vehicle plates from images, checks them against a hotlist database, and plots detection locations on an interactive map. The system uses EasyOCR for text extraction and is deployed as a Streamlit dashboard. On a synthetic dataset of 200 Pakistani-style plates, the system achieved correct plate detection on 80\% of images with a confidence threshold of 0.3. The hotlist alert mechanism provides real-time feedback, and detection history with geolocation data is visualized using Folium maps. The system is designed as a low-cost surveillance tool for urban security applications.
\end{abstract}

\section{Introduction}

License plate recognition (LPR) is widely used in traffic monitoring, toll collection, and law enforcement. Most commercial systems are trained on Western or Chinese plate formats. Pakistani license plates have a distinct format ({2-3 letters}-{3-4 digits}) and varying fonts, which makes existing systems less accurate.

In this paper, we build a complete LPR system from scratch. We generate synthetic Pakistani plate images, use EasyOCR for text extraction, maintain a hotlist database of suspected vehicles, and visualize detections on a map. The goal is a working prototype that can be deployed on a laptop with a webcam.

\section{System Design}

\subsection{Dataset}

We generated 200 synthetic plate images using Python Pillow. Each plate follows the Pakistani format: 2--3 letters, a dash, and 3--4 digits (e.g., \texttt{LEG-456}). Images are 400x120 pixels with a white background, blue borders, and black text. 20 plates (10\%) are randomly hotlisted as "suspicious."

\subsection{Plate Detection and OCR}

We use EasyOCR, a deep learning-based OCR library, to extract text from plate images. EasyOCR is trained on 80+ languages and works well on synthetic text. We filter out common false positives (e.g., \texttt{PAKISTAN}, \texttt{ISLAMABAD}) and only accept detections with confidence above 0.3 and minimum 4 characters.

\subsection{Hotlist Database}

A SQLite database stores two tables: \texttt{hotlist} (plates marked as suspicious, with reason and timestamp) and \texttt{detections} (every detected plate with image name, GPS coordinates, location name, and timestamp). When a scanned plate matches the hotlist, the system triggers an alert.

\subsection{Visualization}

Detection locations are plotted on a Folium map. Hotlisted plates appear as red markers, clean plates as blue. The Streamlit dashboard provides three tabs: scanning, hotlist management, and the tracking map.

\section{Results}

We tested the scanner on all 200 synthetic plate images. Results are summarized below.

\begin{table}[h]
\centering
\caption{Scanner performance on 200 synthetic plates.}
\label{tab:results}
\begin{tabular}{@{}lcc@{}}
\toprule
Metric & Value \\
\midrule
Total plates & 200 \\
Correctly detected & 160 (80.0\%) \\
Missed / incorrect & 40 (20.0\%) \\
Average confidence (correct detections) & 0.87 \\
\bottomrule
\end{tabular}
\end{table}

The system processes each image in under 2 seconds on a CPU. The Streamlit dashboard updates instantly when a plate is logged.

\subsection{Limitations}

Our plates are synthetic -- real-world plates may have varying lighting, angles, and occlusion. The OCR model used (EasyOCR) is general-purpose and was not fine-tuned on plate text. A custom CNN for plate-specific character recognition would improve accuracy.

\section{Conclusion}

We built a functional license plate recognition and hotlist alert system using open-source tools. The system reads plates from images, checks them against a hotlist, and tracks detections on a map. With 80\% accuracy on synthetic data, it demonstrates the feasibility of a low-cost ANPR system for Pakistani plates.

\begin{thebibliography}{5}

\bibitem{easyocr}
A. J\'aid\'ar et al., "EasyOCR: Ready-to-use OCR with 80+ languages," \textit{GitHub repository}, 2022.

\bibitem{anpr}
C. N. E. Anagnostopoulos et al., "License plate recognition from still images and video sequences: A survey," \textit{IEEE Transactions on Intelligent Transportation Systems}, vol. 9, no. 3, pp. 377--391, 2008.

\bibitem{yolo}
J. Redmon et al., "You Only Look Once: Unified, real-time object detection," \textit{Proceedings of CVPR}, pp. 779--788, 2016.

\bibitem{streamlit}
Streamlit Inc., "Streamlit: The fastest way to build and share data apps," \textit{streamlit.io}, 2023.

\bibitem{folium}
Folium Contributors, "Folium: Python library for visualizing geospatial data," \textit{python-visualization.github.io/folium}, 2023.

\end{thebibliography}

\end{document}
